\documentclass[11pt]{article}

\usepackage[margin=2.5cm]{geometry}
\usepackage{amsmath,amssymb,amsthm}
\usepackage{bm}

\title{Exercise Sheet -- Fourier Analysis}
\author{}
\date{}

\begin{document}
\maketitle

This sheet is organized in two parts:
\begin{itemize}
  \item \textbf{Part A -- Fundamental exercises}
  \item \textbf{Part B -- Advanced exercises}
\end{itemize}

The problems are designed to practice: Fourier series (symmetries, basic computations, scaling), Fourier transform in one and several dimensions, core properties (translation, modulation, dilation, differentiation), and some conceptual links with signal processing and machine learning (Fourier features, spectral bias).


\bigskip
\hrule
\bigskip

\section*{Part A -- Fundamental Exercises}

\begin{enumerate}
% -----------------------------
\item \textbf{Even/odd functions and Fourier series structure}

Let \( f(x)=x^2 \) on \( [-\pi,\pi] \), extended as a \(2\pi\)-periodic function.

\begin{enumerate}
\item Determine whether \(f\) is even, odd, or neither.
\item State which Fourier coefficients must be zero (sine or cosine coefficients).
\item Write the general form of the Fourier series of \(f\) (do not compute the integrals).
\end{enumerate}

\bigskip
% -----------------------------
\item \textbf{Fourier series of a square wave}

Consider the \(2\pi\)-periodic function
\[
f(x)=
\begin{cases}
1 & \text{if } 0<x<\pi, \\
-1 & \text{if } -\pi<x<0.
\end{cases}
\]

\begin{enumerate}
\item Show that \(f\) is odd.
\item Deduce that its Fourier series only contains sine terms.
\item Compute the sine coefficients \(b_n\) in the expansion
\[
f(x)\sim \sum_{n=1}^{\infty} b_n \sin(nx).
\]
\end{enumerate}

\bigskip
% -----------------------------
\item \textbf{Scaling property for Fourier series}

Let \(f\) be a \(2\pi\)-periodic function with Fourier series
\[
f(x) \sim \sum_{n=-\infty}^{\infty} c_n e^{inx}.
\]
Define \(g(x)=f(2x)\).

\begin{enumerate}
\item What is the period of \(g\)?
\item Express \(g(x)\) as a Fourier series in the form
\[
g(x) \sim \sum_{k=-\infty}^{\infty} d_k e^{ikx},
\]
and relate the coefficients \(d_k\) to the coefficients \(c_n\).
\item Explain qualitatively what happens to the frequencies when time is compressed by a factor \(2\).
\end{enumerate}

\bigskip
% -----------------------------
\item \textbf{Fourier transform of a Gaussian}

Consider the function
\[
f(x)=e^{-\pi x^2}, \qquad x\in\mathbb{R},
\]
and the Fourier transform
\[
\hat f(\xi)=\int_{-\infty}^{\infty} f(x)\, e^{-2\pi i x\xi}\,dx.
\]

\begin{enumerate}
\item Show that \(\hat f(\xi)=e^{-\pi \xi^2}\) for all \(\xi\in\mathbb{R}\).
\item Comment briefly on the fact that the Gaussian is an eigenfunction of the Fourier transform.
\end{enumerate}

\bigskip
% -----------------------------
\item \textbf{Differentiation in the Fourier transform}

Let \(f\) and its derivative \(f'\) be integrable, and denote by \(\hat f\) and \(\widehat{f'}\) their Fourier transforms.

\begin{enumerate}
\item Prove that
\[
\widehat{f'}(\xi) = 2\pi i \xi\, \hat f(\xi).
\]
You may use integration by parts, assuming suitable decay of \(f\).
\item Explain in words why differentiation in time (or space) amplifies high frequencies.
\end{enumerate}

\bigskip


\end{enumerate}

\bigskip
\hrule
\bigskip

\section*{Part B -- Advanced Exercises}

\begin{enumerate}
\setcounter{enumi}{6} % continue numbering from Part A

% -----------------------------
\item \textbf{Periodic extension of a piecewise linear function}

Let
\[
f(x)=\max(x,0)
\]
defined on \([-\,\pi,\pi]\), and extended as a \(2\pi\)-periodic function outside this interval.

\begin{enumerate}
\item Sketch one period of \(f\).
\item Determine whether \(f\) is even, odd, or neither.
\item Write the formulas defining the Fourier coefficients \(a_n\) and \(b_n\) in the expansion
\[
f(x)\sim \frac{a_0}{2} + \sum_{n=1}^{\infty}\big(a_n\cos(nx)+b_n\sin(nx)\big),
\]
without evaluating the integrals.
\end{enumerate}

\bigskip
% -----------------------------
\item \textbf{Denoising via truncation of a Fourier series}

A signal is modeled as
\[
s(t)=u(t)+\varepsilon(t),
\]
where \(u\) is a smooth underlying signal and \(\varepsilon\) is high-frequency noise.

\begin{enumerate}
\item Explain how truncating the Fourier series of \(s(t)\) (keeping only the low-frequency coefficients) can reduce noise.
\item What are the potential drawbacks of removing too many high-frequency coefficients?
\item Give an example of a situation where high-frequency information is important and should not be discarded.
\end{enumerate}

\bigskip
% -----------------------------
\item \textbf{Convolution and the Fourier transform}

Let \(f\) and \(g\) be integrable functions with Fourier transforms \(\hat f\) and \(\hat g\).
We recall the definition of the convolution \(f*g (x) = \int_{\mathbb{R}} f(z) g (x - z) \mathrm{d}z\).
\begin{enumerate}
\item Show that
\[
\mathcal{F}\{f*g\}(\xi)=\hat f(\xi)\,\hat g(\xi).
\]
\item Explain why this implies that convolution in the time domain corresponds to multiplication in the frequency domain, and how this can be used to design filters.
\end{enumerate}

\bigskip

% -----------------------------
\item \textbf{Convolution and smoothing (basic)}

Let \(g(x)=\frac{1}{2h}\mathbf{1}_{[-h,h]}(x)\), where \(h>0\), and let \(f\) be an integrable function.

\begin{enumerate}
\item Write explicitly the convolution
\[
(f*g)(x) = \int_{-\infty}^{\infty} f(x-y)g(y)\,dy.
\]
\item Interpret this operation in terms of averaging the values of \(f\).
\item Explain intuitively why convolving with \(g\) tends to smooth \(f\).
\end{enumerate}

\bigskip
% -----------------------------
\item \textbf{Fourier transform of an exponential with absolute value}

Consider
\[
f(x)=e^{-a|x|}, \qquad a>0.
\]

\begin{enumerate}
\item Show that \(f\) is an even function.
\item Compute its Fourier transform \(\hat f(\xi)\) explicitly.
\item Describe the decay of \(\hat f(\xi)\) as \(|\xi|\to\infty\), and compare it qualitatively with the Gaussian case.
\end{enumerate}

\bigskip
% -----------------------------
\item \textbf{Two-dimensional Fourier transform of a Gaussian}

Let
\[
f(x,y)=e^{-\pi(x^2+y^2)}, \qquad (x,y)\in\mathbb{R}^2.
\]

\begin{enumerate}
\item Compute the two-dimensional Fourier transform \(\hat f(\xi_1,\xi_2)\).
\item Show that \(\hat f(\xi_1,\xi_2)\) factorizes as the product of two one-dimensional Fourier transforms.
\item Comment briefly on the radial symmetry of both \(f\) and \(\hat f\).
\end{enumerate}

\bigskip
% -----------------------------
\item \textbf{Heat equation solved by Fourier transform}

Consider the one-dimensional heat equation
\[
\partial_t u(t,x)=\partial_{xx}u(t,x), \qquad x\in\mathbb{R},\ t>0,
\]
with initial condition \(u(0,x)=f(x)\).

\begin{enumerate}
\item Apply the Fourier transform in the spatial variable \(x\) to the PDE and write the resulting ordinary differential equation for \(\hat u(t,\xi)\).
\item Solve this ODE for \(\hat u(t,\xi)\) in terms of \(\hat f(\xi)\).
\item Write the solution \(u(t,x)\) using the inverse Fourier transform.
\end{enumerate}

\bigskip
% -----------------------------
\item \textbf{Fourier features}

For \(k\in\mathbb{Z}\), define the real-valued features
\[
\phi_k(x)=\cos(2\pi kx), \qquad \psi_k(x)=\sin(2\pi kx).
\]

\begin{enumerate}
\item Explain how these functions are related to Fourier series.
\item What types of target functions would require many high-frequency features (large \(|k|\)) to be accurately approximated?
\end{enumerate}

\bigskip
% -----------------------------
\item \textbf{Spectral bias of learning algorithms (conceptual)}

Empirically, many learning algorithms tend to learn low-frequency components of a function before high-frequency components.

\begin{enumerate}
\item Reformulate this statement using Fourier analysis language.
\item Give one possible advantage of such a spectral bias.
\item Give one possible drawback of such a spectral bias when learning from data.
\end{enumerate}

\end{enumerate}

\end{document}